\documentclass[10pt, aspectratio=169]{beamer}
\usepackage{../beamer_theme_408}
\title{408考研零基础 --- 计算机网络}
\subtitle{4-9 介质访问控制}
\author{}
\date{}
\begin{document}
\frame{\titlepage}

\begin{frame}{本节目标}
    \begin{enumerate}
        \item 理解信道划分MAC的四种方式：FDM、TDM、WDM、CDM
        \item 掌握随机访问MAC的CSMA/CD协议原理
        \item 了解轮询访问MAC的令牌传递机制
    \end{enumerate}
\end{frame}

\begin{frame}{介质访问控制概述}
    \begin{block}{解决的问题}
        多个站点共享同一传输介质时，如何协调发送以避免冲突？
    \end{block}
    \vspace{0.5em}
    \begin{columns}[t]
        \column{0.33\textwidth}
        \textbf{信道划分MAC}
        \begin{itemize}
            \item FDM
            \item TDM / 统计TDM
            \item WDM
            \item CDM
        \end{itemize}
        \pause
        \column{0.33\textwidth}
        \textbf{随机访问MAC}
        \begin{itemize}
            \item ALOHA
            \item CSMA
            \item CSMA/CD
            \item CSMA/CA
        \end{itemize}
        \pause
        \column{0.33\textwidth}
        \textbf{轮询访问MAC}
        \begin{itemize}
            \item 令牌传递
            \item 无冲突
            \item 负载高时效率高
        \end{itemize}
    \end{columns}
\end{frame}

\begin{frame}{信道划分MAC --- 频分多路FDM}
    \begin{center}
        \begin{tikzpicture}[>=stealth]
            \draw[->] (0,0) -- (8,0) node[right]{频率};
            \fill[red!30] (0,0.2) rectangle (7,0.6);
            \fill[green!30] (0,0.7) rectangle (7,1.1);
            \fill[blue!30] (0,1.2) rectangle (7,1.6);
            \fill[yellow!30] (0,1.7) rectangle (7,2.1);
            \node[right] at (7,0.4) {信道1（频率1）};
            \node[right] at (7,0.9) {信道2（频率2）};
            \node[right] at (7,1.4) {信道3（频率3）};
            \node[right] at (7,1.9) {信道4（频率4）};
            \node[above] at (3.5,2.3) {\textbf{FDM：不同频率同时传输}};
        \end{tikzpicture}
    \end{center}
    \vspace{0.5em}
    \begin{itemize}
        \item 将信道划分为多个不同频段，各站点占用不同频率
        \item 优点：同时传输，互不干扰
        \item 应用：传统电话网络、广播电视
        \item 缺点：频率利用率不高，扩展性差
    \end{itemize}
\end{frame}

\begin{frame}{信道划分MAC --- 时分多路TDM}
    \begin{center}
        \begin{tikzpicture}[>=stealth]
            \draw[->] (0,0) -- (10,0) node[right]{时间};
            \fill[red!30] (0,0.2) rectangle (1.8,0.8);
            \fill[green!30] (1.8,0.2) rectangle (3.6,0.8);
            \fill[blue!30] (3.6,0.2) rectangle (5.4,0.8);
            \fill[yellow!30] (5.4,0.2) rectangle (7.2,0.8);
            \fill[red!30] (7.2,0.2) rectangle (9.0,0.8);
            \fill[green!30] (9.0,0.2) rectangle (10.8,0.8);
            \node[below] at (0.9,0) {A};
            \node[below] at (2.7,0) {B};
            \node[below] at (4.5,0) {C};
            \node[below] at (6.3,0) {D};
            \node[below] at (8.1,0) {A};
            \node[below] at (9.9,0) {B};
            \node[above] at (5,1.2) {\textbf{TDM：不同时间片轮流传输}};
            \draw[decorate,decoration={brace,amplitude=5pt}] (0,1.0) -- (7.2,1.0);
            \node[above] at (3.6,1.5) {一帧(含4个时间片)};
        \end{tikzpicture}
    \end{center}
    \vspace{0.3em}
    \begin{columns}[t]
        \column{0.48\textwidth}
        \textbf{TDM}
        \begin{itemize}
            \item 固定分配时间片
            \item 空闲时片浪费
        \end{itemize}
        \column{0.48\textwidth}
        \textbf{统计TDM}
        \begin{itemize}
            \item 按需分配时间片
            \item 提高信道利用率
            \item 需额外地址信息
        \end{itemize}
    \end{columns}
\end{frame}

\begin{frame}{信道划分MAC --- WDM与CDM}
    \begin{columns}[t]
        \column{0.48\textwidth}
        \textbf{波分多路WDM}
        \begin{itemize}
            \item 不同波长的光信号在同一光纤传输
            \item 本质是光的FDM
            \item 极大提升光纤传输容量
            \item 应用：光纤骨干网
        \end{itemize}
        \vspace{0.5em}
        \begin{center}
            \begin{tikzpicture}[>=stealth]
                \draw[red,->,thick] (0,0) -- (2,0);
                \draw[blue,->,thick] (0,0.3) -- (2,0.3);
                \draw[green,->,thick] (0,0.6) -- (2,0.6);
                \draw[->,thick] (2,0) -- (3.5,0.3);
                \node at (1,-0.3) {多种波长};
                \node at (2.8,-0.3) {光纤};
            \end{tikzpicture}
        \end{center}
        \column{0.48\textwidth}
        \textbf{码分多路CDM}
        \begin{itemize}
            \item 每个站点分配唯一编码序列
            \item 所有站点同时发送，共用同一频率
            \item 接收方用对应编码解码
            \item \textbf{码分多址CDMA}广泛应用（3G）
        \end{itemize}
        \vspace{0.3em}
        \begin{block}{CDMA原理}
            \[
            S \cdot T = 0 \quad (S \neq T)
            \]
            \[
            S \cdot S = 1
            \]
            不同站点的码片序列正交
        \end{block}
    \end{columns}
\end{frame}

\begin{frame}{随机访问MAC --- ALOHA协议}
    \begin{columns}[t]
        \column{0.48\textwidth}
        \textbf{纯ALOHA}
        \begin{itemize}
            \item 任意时刻均可发送
            \item 冲突则随机等待后重发
            \item 最大吞吐率$18\%$
        \end{itemize}
        \vspace{0.5em}
        \begin{center}
            \begin{tikzpicture}[>=stealth]
                \draw[->] (0,0) -- (4,0);
                \fill[blue!30] (0.5,0) rectangle (1.2,0.4);
                \fill[red!30] (1.5,0) rectangle (2.2,0.4);
                \fill[blue!30] (2.8,0) rectangle (3.5,0.4);
                \node[below] at (0.85,-0.2) {站A};
                \node[below] at (1.85,-0.2) {站B};
                \node[below] at (3.15,-0.2) {站A重传};
            \end{tikzpicture}
        \end{center}
        \column{0.48\textwidth}
        \textbf{时隙ALOHA}
        \begin{itemize}
            \item 时间分为等长时隙
            \item 只能在时隙起点发送
            \item 最大吞吐率$37\%$
        \end{itemize}
        \vspace{0.5em}
        \begin{center}
            \begin{tikzpicture}[>=stealth]
                \draw[->] (0,0) -- (4,0);
                \foreach \x in {0,1,2,3} {
                    \draw[dashed] (\x,0) -- (\x,0.5);
                    \node[above] at (\x+0.5,0) {时隙};
                }
                \fill[blue!30] (1,0) rectangle (1.7,0.4);
            \end{tikzpicture}
        \end{center}
    \end{columns}
    \vspace{0.5em}
    \begin{alertblock}{注意}
        ALOHA是最简单的随机访问协议，但冲突概率高，效率低。
    \end{alertblock}
\end{frame}

\begin{frame}{随机访问MAC --- CSMA协议}
    \begin{block}{载波监听多路访问}
        发前先监听信道是否空闲（载波监听），避免盲目发送。
    \end{block}
    \vspace{0.3em}
    \begin{columns}[t]
        \column{0.32\textwidth}
        \textbf{1-持续CSMA}
        \begin{itemize}
            \item 监听信道
            \item 空闲立即发送
            \item 忙则持续监听
        \end{itemize}
        \column{0.32\textwidth}
        \textbf{非持续CSMA}
        \begin{itemize}
            \item 空闲→发送
            \item 忙→随机等待
            \item 减少冲突概率
        \end{itemize}
        \column{0.32\textwidth}
        \textbf{p-持续CSMA}
        \begin{itemize}
            \item 空闲→概率$p$发送
            \item 忙→继续监听
            \item 折中方案
        \end{itemize}
    \end{columns}
    \vspace{0.5em}
    \begin{exampleblock}{考点}
        CSMA仍有冲突（因为传播时延导致监听结果可能过时），于是有了CSMA/CD。
    \end{exampleblock}
\end{frame}

\begin{frame}{CSMA/CD协议 --- 核心原理}
    \begin{center}
        \fbox{\parbox{0.85\textwidth}{\centering \textbf{先听后发 \quad 边听边发 \quad 冲突停发 \quad 随机重发}}}
    \end{center}
    \vspace{0.5em}
    \begin{enumerate}
        \item \textbf{先听后发}：发送前监听信道，空闲才发送
        \item \textbf{边听边发}：发送过程中持续检测冲突
        \item \textbf{冲突停发}：检测到冲突，立即停止发送，发送\textbf{强化干扰信号}(Jam)
        \item \textbf{随机重发}：等待随机时间后重试
    \end{enumerate}
    \vspace{0.5em}
    \begin{block}{截断二进制指数退避算法}
        \[
        t = \text{rand}(0, 2^k - 1) \times 2\tau \quad (k = \min\{\text{重传次数}, 10\})
        \]
        \begin{itemize}
            \item $2\tau$：争用期（$\tau$为最远站点传播时延）
            \item 重传16次仍失败 → 丢弃帧，报告错误
        \end{itemize}
    \end{block}
\end{frame}

\begin{frame}{CSMA/CD --- 最短帧长}
    \begin{block}{为什么需要最短帧长？}
        发送方在发送完帧之前必须检测到冲突，否则认为发送成功。
    \end{block}
    \vspace{0.5em}
    \begin{exampleblock}{最短帧长公式}
        \[
        \text{最短帧长} = 2\tau \times \text{数据传输速率}
        \]
    \end{exampleblock}
    \vspace{0.3em}
    \begin{itemize}
        \item 以太网最短帧长：\textbf{64字节}
        \item $2\tau = 51.2\mu\text{s}$（10Mbps以太网）
        \item 若帧长小于64字节 → 填充至64字节
    \end{itemize}
    \vspace{0.5em}
    \begin{center}
        \begin{tikzpicture}[>=stealth]
            \draw[->] (0,0) -- (6,0);
            \draw[red,thick] (0,0.1) -- (2,0.1);
            \node[above] at (1,0.3) {最短帧发送时间$=2\tau$};
            \node[below] at (1,-0.2) {此期间必须检测到冲突};
        \end{tikzpicture}
    \end{center}
\end{frame}

\begin{frame}{CSMA/CD --- 冲突检测过程}
    \begin{center}
        \begin{tikzpicture}[>=stealth, node distance=1.5cm]
            \node (A) {主机A};
            \node[right=4cm of A] (B) {主机B};
            \draw[->] (A) -- node[above]{发送数据} (B);
            \node[below=0.5cm of A] (t1) {$\tau$时刻：};
            \draw[->] (B) -- node[below]{B也发送} (A);
            \node[below=0.5cm of t1] (t2) {$\tau$时刻：冲突发生};
            \node[below=0.5cm of t2] {最多$2\tau$检测到冲突};
        \end{tikzpicture}
    \end{center}
    \vspace{0.3em}
    \begin{itemize}
        \item A在$t_0$发送，B在$t_0+\tau-\delta$发送
        \item 冲突在$t_0+\tau$发生
        \item A在$t_0+2\tau-\delta$检测到冲突
        \item 结论：最坏情况在$2\tau$后检测到冲突
    \end{itemize}
\end{frame}

\begin{frame}{轮询访问MAC --- 令牌传递}
    \begin{block}{工作原理}
        网络中只有一个令牌（特殊帧）在循环传递，站点拿到令牌才能发送数据。
    \end{block}
    \vspace{0.3em}
    \begin{center}
        \begin{tikzpicture}[>=stealth]
            \foreach \i/\angle/\label in {1/90/A,2/18/B,3/-54/C,4/-126/D,5/-198/E,6/-270/F,7/-342/G} {
                \node[draw,circle,minimum size=0.8cm] (n\i) at (\angle:2cm) {\label};
            }
            \foreach \i in {1,2,3,4,5,6,7} {
                \pgfmathtruncatemacro{\next}{mod(\i,7)+1}
                \draw[->] (n\i) -- (n\next);
            }
            \node[draw,circle,fill=yellow!40,minimum size=0.5cm] at (90:0.5cm) {令牌};
            \draw[->,thick,blue] (90:0.5cm) -- (18:0.5cm);
        \end{tikzpicture}
    \end{center}
    \vspace{0.3em}
    \begin{columns}[t]
        \column{0.48\textwidth}
        \textbf{优点}
        \begin{itemize}
            \item 无冲突
            \item 负载高时效率高
            \item 可设置优先级
        \end{itemize}
        \column{0.48\textwidth}
        \textbf{缺点}
        \begin{itemize}
            \item 令牌维护复杂
            \item 单点故障（令牌丢）
            \item 轻负载时效率低
        \end{itemize}
    \end{columns}
\end{frame}

\begin{frame}{三种MAC方式对比}
    \begin{table}
        \centering
        \begin{tabular}{|c|c|c|c|}
            \hline
            \textbf{特性} & \textbf{信道划分} & \textbf{随机访问} & \textbf{轮询} \\
            \hline
            冲突 & 无 & 有 & 无 \\
            轻负载效率 & 低 & 高 & 低 \\
            重负载效率 & 高 & 低（冲突多） & 高 \\
            公平性 & 好 & 较差 & 好 \\
            典型协议 & FDM/TDM/CDMA & CSMA/CD & 令牌环 \\
            \hline
        \end{tabular}
    \end{table}
    \vspace{0.5em}
    \begin{itemize}
        \item 信道划分：预先分配资源，适合恒定流量
        \item 随机访问：按需竞争，适合突发流量
        \item 轮询访问：有序访问，适合负载较重场景
    \end{itemize}
\end{frame}

\begin{frame}{本节总结}
    \begin{enumerate}
        \item 信道划分MAC：FDM（频率）、TDM（时间片）、WDM（波长）、CDM（编码）
        \item 随机访问MAC：ALOHA（纯18\%/时隙37\%）、CSMA（监听后发）、CSMA/CD（先听后发、边听边发、冲突停发、随机重发），最短帧长64B
        \item 轮询访问MAC：令牌传递，无冲突，重负载效率高
    \end{enumerate}
    \vspace{1em}
    \begin{center}
        \textcolor{blue}{\textbf{下节预告：4-10 网络层功能与IP}}
    \end{center}
\end{frame}
\end{document}
